\pdfoutput=1
% Options for packages loaded elsewhere
\PassOptionsToPackage{unicode}{hyperref}
\PassOptionsToPackage{hyphens}{url}
\documentclass[
  11pt,
]{article}
\usepackage{xcolor}
\usepackage[margin=1in]{geometry}
\usepackage{amsmath,amssymb}
\setcounter{secnumdepth}{-\maxdimen} % remove section numbering
\usepackage{iftex}
\ifPDFTeX
  \usepackage[T1]{fontenc}
  \usepackage[utf8]{inputenc}
  \usepackage{textcomp} % provide euro and other symbols
\else % if luatex or xetex
  \usepackage{unicode-math} % this also loads fontspec
  \defaultfontfeatures{Scale=MatchLowercase}
  \defaultfontfeatures[\rmfamily]{Ligatures=TeX,Scale=1}
\fi
\usepackage{lmodern}
\ifPDFTeX\else
  % xetex/luatex font selection
\fi
% Use upquote if available, for straight quotes in verbatim environments
\IfFileExists{upquote.sty}{\usepackage{upquote}}{}
\IfFileExists{microtype.sty}{% use microtype if available
  \usepackage[]{microtype}
  \UseMicrotypeSet[protrusion]{basicmath} % disable protrusion for tt fonts
}{}
\makeatletter
\@ifundefined{KOMAClassName}{% if non-KOMA class
  \IfFileExists{parskip.sty}{%
    \usepackage{parskip}
  }{% else
    \setlength{\parindent}{0pt}
    \setlength{\parskip}{6pt plus 2pt minus 1pt}}
}{% if KOMA class
  \KOMAoptions{parskip=half}}
\makeatother
\setlength{\emergencystretch}{3em} % prevent overfull lines
\providecommand{\tightlist}{%
  \setlength{\itemsep}{0pt}\setlength{\parskip}{0pt}}
\usepackage[]{natbib}
\bibliographystyle{plainnat}
\usepackage{bookmark}
\IfFileExists{xurl.sty}{\usepackage{xurl}}{} % add URL line breaks if available
\urlstyle{same}
\hypersetup{
  pdftitle={Example 2 - Software Developer},
  hidelinks,
  pdfcreator={LaTeX via pandoc}}

\title{Example 2 - Software Developer}
\author{}
\date{May 2025}

\begin{document}
\maketitle

\subsection{Software Developer}\label{software-developer}

\textbf{1. Job Composition}

\begin{itemize}
\tightlist
\item
  Thinking: 60\% (coding, problem-solving, design)
\item
  Physical: 5\% (typing, basic computer interaction)
\item
  Social: 20\% (stakeholder communication, requirements gathering)
\item
  Perception: 10\% (code review, testing)
\item
  Adaptation: 5\% (handling changing requirements)
\end{itemize}

\textbf{2. Human Advantage Score}

\begin{itemize}
\tightlist
\item
  Thinking: 4/10 (good at code generation but struggles with
  larger-scale planning)
\item
  Physical: 1/10 (AI doesn't need physical interaction to code)
\item
  Social: 5/10 (struggles with nuanced stakeholder management)
\item
  Perception: 4/10 (effective at code analysis)
\item
  Adaptation: 4/10 (can adapt to changes but misses context)
\end{itemize}

\textbf{3. Job Resilience Score}

\begin{itemize}
\tightlist
\item
  Total: 4.05/10
\item
  Beginner: 3.05 (4.05-1)
\item
  Advanced: 4.05
\item
  Master: 5.05 (4.05+1)
\end{itemize}

\textbf{4. Result:}

\begin{itemize}
\tightlist
\item
  Beginner: High Risk (1-5 years)
\item
  Advanced: Medium Risk (5-10 years)
\item
  Master: Medium Risk (5-10 years)
\end{itemize}

\textbf{5. Adoption Factors Check: }High salaries justify AI
replacement, potentially moving advanced developers toward High Risk,
while master developers with stronger stakeholder relationships remain
in Medium Risk\textbf{.}

\end{document}
